\section{An outer game with affine hints}
\label{sec:outer}

This Section constructs the outer game and proves the two properties
used in Section~\ref{sec:composition}: soundness with shared affine hints,
and approximate joint uniformity of pulled-back affine functions.

Smooth equation--variable sampling and Fourier decoding have antecedents
in Khot and Safra~\cite[Sections~3.2--3.4 and~5.2]{KS13}.
Khot, Minzer and Safra~\cite[Sections~3.3--3.4]{KMS17} use shared linear
information and zero-advice coordinates; Dinur et al.~\cite[Section~5.2,
Lemma~5.4]{DKKMS2025} give a related shared-subspace formulation.
We prove both properties for the exact questions and hints defined below.

Let $\mathcal E=(E^{(1)},\ldots,E^{(M_0)})$ be the nonempty uniform
equation list supplied by Theorem~\ref{thm:max3lin}. An equation question
$E$ consists of its ordered three variable names $v(E,1),v(E,2),v(E,3)$
and its right-hand side $\tau_E\in\F_2$. Thus it asks that
\[
 x_{v(E,1)}+x_{v(E,2)}+x_{v(E,3)}=\tau_E.
\]
The list may have repeated equations, and variable names may repeat within
an equation. Sampling is uniform over list entries, but equal equation
questions have identical encodings: an index distinguishing equal list
entries is not part of a player's question.

\subsection{Questions, fresh slots, and canonical functions}

For an equation $E$, put
\[
 A(E)=\{a\in\F_2^3:
        a_1+a_2+a_3=\tau_E,\quad
        a_p=a_{p'}\text{ whenever }v(E,p)=v(E,p')\}.
\]
This is a nonempty affine space of dimension at most two. Indeed, after
identifying repeated variable names, there are $k\le3$ free bit slots and
one nonzero linear equation: at least one multiplicity is odd because the
sum of the multiplicities is three. The equation therefore has rank one,
and $\dim A(E)=k-1$.

Fix integers $t\ge t_0\ge1$. An \emph{outer sample} $e$ is drawn as follows.
First draw $t$ independent uniform list entries, with resulting equation
questions $E_1,\ldots,E_t$, and let
\[
 U=(E_1,\ldots,E_t),\qquad A_U=\prod_{j=1}^t A(E_j).
\]
Next draw a uniform $t_0$-element subset $H\subseteq[t]$. For each $j\in H$,
independently draw $p_j\in[3]$ uniformly. The second question is the ordered,
tagged tuple $V=(V_1,\ldots,V_t)$, where
\[
 V_j=\begin{cases}
       (\mathrm{variable},v(E_j,p_j)),&j\in H,\\
       (\mathrm{equation},E_j),&j\notin H.
     \end{cases}
\]
In particular, a variable question contains its variable name and the
coordinate $j$ in the ordered tuple, but neither its position $p_j$ nor
its originating equation. Set
\[
 A_V=\prod_{j=1}^t A_{V,j},\qquad
 A_{V,j}=\begin{cases}\F_2,&j\in H,\\\F_2^3,&j\notin H.\end{cases}
\]
Every displayed slot in this product is unrestricted. Even when the same
variable name occurs more than once, no equalities between these slots are
imposed. Likewise, the slots belonging to distinct coordinates of $A_U$
are separate; its only equalities between repeated names are those inside
a single factor $A(E_j)$. These fresh slots are essential both to independent
repetition and to the distribution of affine forms below.

Define the affine map $\pi_e:A_U\to A_V$ coordinatewise by
\begin{equation}\label{eq:outer-projection}
 (\pi_e(a))_j=\begin{cases}(a_j)_{p_j},&j\in H,\\a_j,&j\notin H.\end{cases}
\end{equation}
The verifier accepts labels $a\in A_U$ and $b\in A_V$ precisely when
$\pi_e(a)=b$. The map need not be onto: the unrestricted full triples on
the second side need not satisfy their equations. Pullback is the linear
map
\[
 \pi_e^*:D(A_V)\longrightarrow D(A_U),\qquad L\longmapsto L\circ\pi_e,
\]
and preserves the constant-one function.

We specify exactly what a question and an affine hint reveal. Every actual
$U$ or $V$ is the tuple just defined, without extra information identifying
its occurrence in $e$. Put its alphabet in lexicographic order as a subset
of its displayed bit slots. A function on that alphabet is represented by
its complete truth table in this order. In particular, a hint in $D(A_U)$
or $D(A_V)$ is the function itself, with this canonical encoding. It is not
accompanied by a formula, its original coefficients on another alphabet,
or an encoding of the map through which it was pulled back. Equal functions
on the same actual question have equal encodings. The same convention will
apply to Boolean queries in the bit proofs. Thus a response function at a
given question cannot distinguish different descriptions of the same query.

For an integer $q\ge0$, the \emph{game with $q$ affine hints} additionally
draws independent uniform functions $L_1,\ldots,L_q\in D(A_V)$. The first
player receives
\[
 (U,\pi_e^*L_1,\ldots,\pi_e^*L_q)
\]
and the second receives $(V,L_1,\ldots,L_q)$. The acceptance condition is
still $\pi_e(a)=b$. Giving independent shared or private randomness to these
players does not change the upper bounds on its value.

\subsection{Soundness despite the affine hints}

Consider first the ordinary equation-versus-variable game: draw a uniform
list entry $E$, choose a uniform position $p\in[3]$, give $E$ to the first
player and only $v(E,p)$ to the second, and ask for $a\in A(E)$ and a bit
$b$. Accept when $a_p=b$. Equivalently the first answer alphabet can be
$\F_2^3$ with invalid answers rejected, so the answer sizes are at most
eight and two.

\begin{lemma}\label{lem:base-equation-game}
If no assignment satisfies more than $3/4$ of the equation list, the
ordinary equation-versus-variable game has value at most $11/12$.
Consequently its ordinary $r$-fold repetition has value at most $c^r$, for
the effective universal constant $c<1$ of
Theorem~\ref{thm:parallel-repetition}, for every $r\ge0$.
\end{lemma}
\begin{proof}
Fix a deterministic second strategy. Its bit for each variable name defines
a global assignment. On at least $1/4$ of the sampled equations this
assignment violates the equation. Every valid first answer to one of these
equations disagrees with this assignment in at least one of the three
positions, so a uniform position detects disagreement with probability at
least $1/3$. An invalid answer is rejected in every position. The overall
failure probability is therefore at least $1/12$. Repeated variable names
do not affect the argument: equality at all three positions would still
give precisely the violated assignment on that equation. The repeated
bound follows from Theorem~\ref{thm:parallel-repetition}; its hypotheses
depend only on this value bound and the two answer sizes.
\end{proof}

\begin{lemma}[Hint soundness]\label{lem:hint-soundness}
Suppose no assignment satisfies more than $3/4$ of the equation list. For
every $q\ge0$ and $t\ge t_0\ge1$, the value of the game with $q$ affine
hints is at most
\begin{equation}\label{eq:hint-bound}
 \bigl(1-2^{-q}(1-c)\bigr)^{t_0}.
\end{equation}
The bound also applies to strategies using arbitrary independent shared
or private randomness. In particular, for fixed $q$ and any prescribed
rational $s>0$, an effective choice of $t_0$ makes the bound strictly less
than $s$, simultaneously for every $t\ge t_0$, every equation-list size,
and every completeness parameter in Theorem~\ref{thm:max3lin}.
\end{lemma}
\begin{proof}
Use the unrestricted slots of $A_V$ to generate the hints. Let $z_j$ denote
the one-bit coordinate when $j\in H$, and let $z_{jp}$, $p\in[3]$, denote
the three coordinates when $j\notin H$. A uniform affine form has the
unique expression
\begin{equation}\label{eq:hint-coefficients}
 L_\ell(z)=\kappa_\ell+
       \sum_{j\in H}\beta_{\ell j}z_j+
       \sum_{j\notin H}\sum_{p=1}^3\beta_{\ell jp}z_{jp}.
\end{equation}
All coefficients and constants in these expressions, for $1\le\ell\le q$,
are independent uniform bits. They can be sampled after $H$ and before
any equation or position, because the abstract slots in this expression
depend only on $H$. Thus their distribution is independent of all sampled
equations and positions.

Define the active coordinates by
\[
 A=\{j\in H:\beta_{\ell j}=0\text{ for every }1\le\ell\le q\}.
\]
For each $j\in H$ its vector of $q$ coefficients is uniform in $\F_2^q$,
independently across $j$. Consequently
\begin{equation}\label{eq:active-binomial}
 |A|\sim\operatorname{Bin}(t_0,2^{-q}).
\end{equation}
Conditional on $H$, the active set and the coefficient bits are independent
of the equation and position samples. The equations are also independent
of $H$. When $q=0$, all hidden coordinates are active and
\eqref{eq:active-binomial} means $|A|=t_0$ deterministically.

For the analysis give both players the following additional information:
$H$, every coefficient and constant in \eqref{eq:hint-coefficients}, and
all equation samples and all chosen positions outside $A$. Denote this
background by $W$. It determines $A$ without inspecting its equation or
position samples. Conditional on any possible $W$, the pairs
\[
 (E_j,p_j),\qquad j\in A,
\]
are still independent samples of a uniform equation-list entry and an
independent uniform position in $[3]$. This follows directly from the
product sampling distribution: once $H$ and the coefficients are fixed,
$A$ is fixed, and conditioning on samples in its complement does not
change the samples in $A$.

We now show that, for fixed $W$, ordinary repeated-game questions in the
active coordinates simulate each player's original information exactly.
List $A$ in increasing order, so that its $|A|$ coordinates can be identified
with an ordinary repetition.
\begin{enumerate}[label=\textup{(\roman*)}]
\item The first player receives the active equation questions $(E_j)_{j\in A}$.
Together with $W$, these reconstruct all of $U$ and hence $A_U$.
For every hint the active summands in \eqref{eq:hint-coefficients} are
identically zero. On any $a\in A_U$, therefore,
\[
 (\pi_e^*L_\ell)(a)=\kappa_\ell+
    \sum_{j\in H\setminus A}\beta_{\ell j}(a_j)_{p_j}+
    \sum_{j\notin H}\sum_{p=1}^3\beta_{\ell jp}(a_j)_p.
\]
Every position appearing on the right is in the background. Thus the first
player can compute the canonical truth table of each pulled-back hint,
without receiving any active $p_j$ or active variable question from the
second side. Although the domain $A_U$ depends on active equations, these
are exactly the first player's own repeated-game questions.
\item The second player receives only the active variable names
$(v(E_j,p_j))_{j\in A}$. Together with $W$ and the tags specified by $H$,
these reconstruct $V$ and $A_V$. Formula~\eqref{eq:hint-coefficients}, whose
active coefficients vanish, computes every direct hint on this alphabet.
Thus the second player can compute its canonical truth table without
receiving an active equation or the chosen position within one.
\end{enumerate}
These reconstructions also explain why the representation of a hint
matters. No metadata attached to a zero summand is present in the actual
message. Restricting the zero function to an affine domain cannot create
an offset or convey a hidden position. All nonzero contributions and their
possible affine offsets are determined by the background and the receiving
player's own question.

Apply the original response functions to the reconstructed information,
and restrict their answers to the active coordinates. This defines a
strategy for the ordinary $|A|$-fold repeated game, with $W$ fixed as a
parameter. Acceptance of the full outer sample implies acceptance in
every active coordinate, since \eqref{eq:outer-projection} imposes exactly
the ordinary equation-versus-variable comparison there. By
Lemma~\ref{lem:base-equation-game}, its conditional probability is at most
$c^{|A|}$; for $A=\varnothing$ this is the trivial bound one.

This reasoning includes coincidences between variable names or equation
questions in different coordinates, and coincidences with the exposed
background. Independent samples may coincide without ceasing to have a
product distribution. Once $W$ is fixed, the simulated strategies may use
all information about such coincidences obtainable from their respective
questions, just as arbitrary strategies for ordinary repetition may do.
No equality between fresh answer slots, nor any further information about
an active question on the opposite side, has been supplied.

Average over $W$ and use \eqref{eq:active-binomial}. The binomial generating
function gives
\[
 \Prob[\text{outer acceptance}]
 \le\E c^{|A|}
 =\sum_{r=0}^{t_0}\binom{t_0}{r}(2^{-q}c)^r
                       (1-2^{-q})^{t_0-r}
 =\bigl(1-2^{-q}(1-c)\bigr)^{t_0}.
\]
One may first fix all independent random seeds of the players, apply the
bound, and then average; this covers the stated random strategies.

Finally replace $c$, if necessary, by an effectively chosen rational
number $\bar c$ with $c\le\bar c<1$. Successively multiplying the rational
number $1-2^{-q}(1-\bar c)<1$ finds an integer $t_0$ whose power is strictly
less than $s$. This finite procedure is independent of $t$, the instance,
and its completeness error. It supplies the asserted effective choice.
\end{proof}

\subsection{Joint marginal smoothness}

The following estimate concerns the joint distribution of several forms
pulled back through the same outer sample. Conditioning is on $U$ only;
in particular it is not conditioning on any previously received hint.
For finite probability distributions $\mu,\nu$, use
$\TV(\mu,\nu)=\frac12\sum_x\abs{\mu(x)-\nu(x)}$.

\begin{lemma}[Joint smoothness]\label{lem:smoothness}
Fix any actual first question $U$, integers $d\ge1$ and $t\ge t_0\ge1$,
and draw the remaining outer sample conditional on $U$. Draw
$L_1,\ldots,L_d$ independently and uniformly from $D(A_V)$. Then
\begin{equation}\label{eq:smoothness-bound}
 \TV\left(\mathcal L\bigl((\pi_e^*L_1,\ldots,\pi_e^*L_d)\mid U\bigr),
                 \operatorname{Unif}(D(A_U)^d)\right)
 \le\frac12\sqrt{(2^{4dt_0}-1)\frac{t_0^2}{t}}.
\end{equation}
The bound is uniform in $U$ and the equation list, including repeated
variable names and zero-dimensional local alphabets. For fixed $d,t_0$
and positive rational $\Delta$, an effective integer $t\ge t_0$ makes
the right-hand side strictly less than $\Delta$.
\end{lemma}
\begin{proof}
Write $A_j=A(E_j)$ and $r_j=\dim A_j\le2$. Choose, using $E_j$ alone,
an affine origin $o_j\in A_j$ and an injective linear map
$T_j:\F_2^{r_j}\to\F_2^3$ with image the translation space of $A_j$.
Thus every label has a unique expression
\[
 a_j=o_j+T_ju_j,\qquad u_j\in\F_2^{r_j}.
\]
For example, lexicographic search and binary Gaussian elimination make
these choices canonical and effective. In these product coordinates a
form on $A_U$ consists of one constant bit and a linear part in each
coordinate group $\F_2^{r_j}$. A $d$-tuple of forms consequently consists
of $d$ constant bits and groups
\[
 \xi_j\in(\F_2^{r_j})^d,\qquad 1\le j\le t.
\]

Generate the $d$ forms by their independent raw coefficients and constants
on the unrestricted slots of $A_V$, as in
\eqref{eq:hint-coefficients}. After pullback, their constant bits are the
original independent uniform constants plus offsets determined by the
other coefficients, the positions, and the origins $o_j$. Conditional
on all those other data, this addition is a translation of $\F_2^d$.
Hence the resulting constants are jointly uniform and independent of
every linear part, of $H$, and of the chosen positions. Their common
uniform factor contributes no total-variation distance, so it suffices
to analyze the joint linear parts $(\xi_1,\ldots,\xi_t)$.

Fix $H$. If $j\notin H$, a single form has independent uniform coefficient
vector $w\in\F_2^3$ on its three raw slots; its pulled-back linear part
is $T_j^{\mathsf T}w$. The map $T_j^{\mathsf T}$ has rank $r_j$ and is
onto $\F_2^{r_j}$. Every fiber therefore has the same size, so its output
is uniform. All $d$ forms in this group are independent, and their joint
law is the uniform law $\mu_j$ on $(\F_2^{r_j})^d$.

If $j\in H$, condition first on its chosen position $p_j$. The $\ell$th
form has a uniform bit coefficient $\beta_{\ell j}$ on that slot, so its
linear part is $\beta_{\ell j}v_{j,p_j}$, where $v_{j,p}$ is the linear
part of the $p$th coordinate of $o_j+T_ju_j$. Average over the single
uniform choice of $p_j$ shared by all $d$ forms, to obtain a distribution
$Q_j$ on $(\F_2^{r_j})^d$. Its density $\rho_j$ with respect to $\mu_j$
satisfies
\begin{equation}\label{eq:local-density}
 \E_{\mu_j}\rho_j=1,\qquad
 0\le\rho_j\le2^{r_jd}\le2^{2d}.
\end{equation}
Indeed, every point of this finite space has uniform mass $2^{-r_jd}$,
whereas its $Q_j$ mass is at most one. This argument allows a position's
linear part to be zero or to coincide with another position's linear
part. If $r_j=0$, the space is a singleton and $\rho_j=1$.

For fixed $H$, the linear-part groups are independent across $j$:
their coefficients are independent, and the positions in distinct hidden
coordinates are independently sampled. Fresh slots give this independence
even when variable names in different equations coincide. Relative to
the product uniform measure $\mu=\prod_{j=1}^t\mu_j$, their joint density
is therefore
\[
 g_H(\xi_1,\ldots,\xi_t)=\prod_{j\in H}\rho_j(\xi_j).
\]
The densities $\rho_j$ depend on $U,d,j$ but not on the other elements
of $H$. Since $H$ is uniform and independent of $U$, the desired mixed
density is $g=\E_H g_H$.

Let $H'$ be an independent uniform $t_0$-element subset of $[t]$.
If $H\cap H'=\varnothing$, independence of the coordinate groups and
\eqref{eq:local-density} give
$\E_\mu g_Hg_{H'}=1$. For every $H,H'$, the pointwise bound
$g_Hg_{H'}\le2^{4dt_0}$ gives the sufficient bound
$\E_\mu g_Hg_{H'}\le2^{4dt_0}$. Moreover, for fixed $H$,
\[
 \Prob[H\cap H'\ne\varnothing]
 \le\sum_{j\in H}\Prob[j\in H']
 =\frac{t_0^2}{t}.
\]
Since $\E_\mu g=1$, it follows that
\begin{align*}
 \E_\mu(g-1)^2
 &=\E_{H,H'}\E_\mu g_Hg_{H'}-1\\
 &\le(2^{4dt_0}-1)\Prob[H\cap H'\ne\varnothing]\\
 &\le(2^{4dt_0}-1)\frac{t_0^2}{t}.
\end{align*}
Cauchy--Schwarz now yields
$\TV(g\mu,\mu)=\frac12\E_\mu\abs{g-1}
\le\frac12\sqrt{\E_\mu(g-1)^2}$, proving
\eqref{eq:smoothness-bound}, including the independent uniform constants.

All estimates used only $r_j\le2$, so the bound is uniform as asserted.
For effectiveness, it is enough to choose an integer
\[
 t\ge t_0,\qquad
 t>\frac{(2^{4dt_0}-1)t_0^2}{4\Delta^2}.
\]
The right-hand side is an explicitly computable rational number for fixed
$d,t_0,\Delta$. No parameter depending on $U$ occurs in this choice.
\end{proof}

\begin{remark}\label{rem:joint-smoothness-use}
Because \eqref{eq:smoothness-bound} is uniform in $U$, it may be averaged
over any distribution of $U$ and any extra randomness independent of the
outer sample. In particular, if a bounded statistic is chosen as a function
of $U$ and some of the pulled-back forms, one may apply Lemma~\ref{lem:smoothness} to the
joint collection of all forms used by that statistic. For a statistic
bounded in absolute value by $B$, the change in expectation is at most
$2B$ times the displayed total-variation bound. This does not assert
uniformity of one form conditional on previously revealed forms.
\end{remark}

Finally, these outer games are explicit finite objects. Their complete
sampling space has
\[
 M_0^t\binom{t}{t_0}3^{t_0}
\]
equally likely draws, counting repeated list entries with their
multiplicities. There are at most $2^{2t}$ labels at a first question and
$2^{3t-2t_0}$ labels at a second question. Their affine spaces, functions,
canonical truth tables, and projections are all enumerable by binary
linear algebra and finite evaluation. For fixed $t,t_0,q,d$ these local
enumerations have constant size independent of the input. Comparing
actual question encodings and identifying equal ones takes polynomial
bit complexity. Thus the fresh-slot and canonical-function conventions
introduce no oracle or uncharged lookup table.
